\documentclass[10pt,twocolumn]{article}
\usepackage[T1]{fontenc}
\usepackage[utf8]{inputenc}
\usepackage{lmodern}
\usepackage[margin=1.8cm,top=2.4cm,bottom=2cm,columnsep=0.6cm]{geometry}
\usepackage{fancyhdr}
\usepackage{titlesec}
\usepackage{booktabs}
\usepackage{amsmath,amssymb,amsfonts}
\usepackage{microtype}
\usepackage{xcolor}
\usepackage{mdframed}
\usepackage{parskip}
\usepackage{enumitem}
\usepackage{hyperref}

\definecolor{rulered}{RGB}{180,30,30}
\definecolor{boxgray}{RGB}{245,245,245}
\definecolor{keywordgray}{RGB}{100,100,100}

\pagestyle{fancy}
\fancyhf{}
\fancyhead[L]{\textsc{\small Claude Mag}}
\fancyhead[C]{\small\textit{Machine Learning Research Digest}}
\fancyhead[R]{\small 2026-09-07}
\fancyfoot[C]{\small\thepage}
\renewcommand{\headrulewidth}{0.8pt}

\titleformat{\section}{\normalsize\bfseries\color{rulered}}{}{0pt}{}%
  [\vspace{-4pt}\color{rulered}\rule{\columnwidth}{0.4pt}]
\titlespacing{\section}{0pt}{8pt}{4pt}

\hypersetup{colorlinks=true,urlcolor=rulered,linkcolor=black}

\begin{document}
\setlength{\parindent}{0pt}
\setlength{\parskip}{4pt}

{\small\color{keywordgray}\textbf{Keywords:}
  unified multimodal models \textbullet{}
  visual understanding \textbullet{}
  image generation \textbullet{}
  task-decoupled architecture}\par\vspace{4pt}

{\LARGE\bfseries\raggedright
  Uncovering Understanding--Generation Synergy
  in Unified Multimodal Models}\par\vspace{4pt}

{\large\itshape\raggedright
  Joint Training Helps Both Tasks --- But Only
  When Visual Computation Is Architecturally
  Decoupled by Objective}\par\vspace{6pt}

{\small\textbf{By} Penghao Wu, Haiwen Diao, Weichen Fan, Lewei Lu,
  Dahua Lin, Ziwei Liu
  (NTU S-Lab; SenseTime Research)}\par\vspace{2pt}

{\small\color{keywordgray}arXiv:2609.01607 \textbullet{} Preprint, September 2026}
\par\vspace{2pt}\noindent{\color{rulered}\rule{\columnwidth}{0.6pt}}\vspace{6pt}

Unified multimodal models (UMMs) jointly perform visual understanding
and image generation within a single model. Do the two objectives help
each other, compete, or merely coexist? This paper answers through a
three-level investigation --- representation, task, and system --- in a
controlled setting without pretrained vision priors. The verdict: synergy
is real at the representation level, but forces both objectives through
the same computation path and one will dominate. A task-decoupled
architecture that specialises conflicting visual computation while
preserving shared semantic interaction avoids this failure mode and
yields consistent gains on both objectives simultaneously.

\begin{mdframed}[backgroundcolor=boxgray,linecolor=rulered,linewidth=1pt,
    innerleftmargin=6pt,innerrightmargin=6pt,
    innertopmargin=6pt,innerbottommargin=6pt]
\textbf{\small AT A GLANCE}\par\vspace{4pt}
\begin{description}[leftmargin=0pt,labelsep=4pt,itemsep=2pt,parsep=0pt,
    font=\small\bfseries]
  \item[Problem:] Does joint understanding/generation training help,
    hurt, or have no effect when both objectives share the same model?
  \item[Why it matters:] Unified models promise efficiency, but training
    instability and objective conflict limit their practical quality.
  \item[Method:] Three-level analysis (representation, task, system)
    with task-decoupled architectural fix.
  \item[Result:] Decoupled architecture improves understanding by 3.2\%
    and generation FID by 21\% vs.\ naive joint baseline.
  \item[Evidence:] Controlled experiments at 300M--3B parameters; VQA,
    image captioning, and text-to-image generation benchmarks.
\end{description}
\end{mdframed}

\section{The Big Picture}

Large vision-language models have converged on two paradigms: dedicated
understanding models (e.g.\ LLaVA, InternVL) and dedicated generation
models (e.g.\ Stable Diffusion, DALL-E). Unifying both in a single model
is appealing: shared representations could reinforce each other, reduce
model count in deployment, and enable joint reasoning over perception
and synthesis.

Yet practitioners have observed that naive joint training frequently
degrades one or both tasks compared to separate specialist models. The
reasons are poorly understood. This paper systematically characterises
the conflict and proposes a structural solution.

\section{Why Existing Approaches Fall Short}

Existing unified models (e.g.\ Emu, Show-o, Janus) attempt joint training
but either (a) use a pretrained vision encoder that introduces its own
inductive bias and masks the training interaction, or (b) report
mixed results with no theoretical explanation. Three failure modes emerge:

\textbf{Objective dominance.} When both tasks share the same visual
encoding layers, the task with stronger or more numerous gradients
suppresses the other --- whichever has lower loss tends to pull
parameters away from the other objective's optimal region.

\textbf{Representation conflict.} Optimal representations for understanding
encode semantic categories and attributes; optimal representations for
generation encode low-level textures and pixel statistics. Forcing both
through the same encoder causes a compromise that is suboptimal for both.

\textbf{Confounded benchmarks.} Papers using pretrained encoders cannot
isolate the effect of joint training from the encoder's pretrained features.

\section{Core Mechanism}

The paper identifies that conflict arises in \emph{early visual encoding
components} (tokeniser, first few attention layers) but not in \emph{semantic
interaction layers} (cross-attention, late transformer blocks).

This motivates the task-decoupled architecture:
\begin{itemize}[itemsep=2pt,parsep=0pt]
  \item \textbf{Specialised early encoders:} separate but lightweight
    visual tokeniser branches for understanding ($\mathrm{Enc}_u$) and
    generation ($\mathrm{Enc}_g$).
  \item \textbf{Shared semantic backbone:} a single transformer stack
    processes fused semantic representations from both branches.
\end{itemize}

Decoupling allows each encoder to optimise its own visual features while
the shared backbone captures cross-modal semantics that benefit both tasks.

\section{Technical Formulation}

\subsection{Joint Training Objective}

The standard (conflict-prone) joint objective is:
\[
  \mathcal{L}_\mathrm{joint}
    = \lambda_u \mathcal{L}_\mathrm{understanding}
    + \lambda_g \mathcal{L}_\mathrm{generation}
\]
where $\mathcal{L}_u$ is cross-entropy over VQA/captioning tokens and
$\mathcal{L}_g$ is a denoising diffusion loss or autoregressive token loss.

\subsection{Dominance Score}

To quantify asymmetric degradation, the dominance score $D$ is:
\[
  D = \frac{\mathrm{Acc}_u^\mathrm{joint} / \mathrm{Acc}_u^\mathrm{separate}}
           {\mathrm{Acc}_g^\mathrm{joint} / \mathrm{Acc}_g^\mathrm{separate}}
\]
$D \gg 1$: understanding dominates. $D \ll 1$: generation dominates.
Naive joint training produces $D \neq 1$, confirming dominance. The
task-decoupled model produces $D \approx 1$.

\subsection{Decoupled Forward Pass}

Given image $x$ and text tokens $t$:
\begin{align}
  z_u &= \mathrm{Enc}_u(x) \\
  z_g &= \mathrm{Enc}_g(x) \\
  z_\mathrm{shared} &= \mathrm{TransformerShared}(z_u, z_g, t) \\
  \mathcal{L}_\mathrm{decoupled} &= \mathcal{L}_u(z_\mathrm{shared})
    + \mathcal{L}_g(z_\mathrm{shared})
\end{align}

Both encoders see the shared backbone's gradients, enabling beneficial
cross-task signal, while their independent parameters specialise without
conflict.

\subsection{Representation-Level Synergy (CKA)}

Synergy is measured via Centred Kernel Alignment (CKA) between
understanding-task-optimal and generation-task-optimal representations.
After decoupled training, CKA of the encoder outputs is lower (more
specialised), but CKA of the semantic backbone outputs is higher (more
shared) --- confirming specialisation at the encoder and consolidation
at the semantic level.

\section{Learning or Inference Procedure}

Training follows standard autoregressive pretraining on interleaved
image-text data, with the following modification:
\begin{enumerate}[itemsep=2pt,parsep=0pt]
  \item Separate learnable visual tokenisers are initialised from
    a shared initialisation but trained independently.
  \item Gradients from $\mathcal{L}_u$ update $\mathrm{Enc}_u$ and
    the shared backbone; gradients from $\mathcal{L}_g$ update
    $\mathrm{Enc}_g$ and the shared backbone.
  \item The shared backbone receives signal from both tasks at every step.
\end{enumerate}

At inference, only the relevant encoder branch is activated: understanding
queries use $\mathrm{Enc}_u$; generation queries use $\mathrm{Enc}_g$.

\section{What the Guarantee Says}

The paper proves that under mild smoothness conditions on the loss
landscape, the decoupled architecture has a strictly smaller gradient
conflict norm than the fully shared architecture:
\[
  \|\nabla_\theta \mathcal{L}_u + \nabla_\theta \mathcal{L}_g\|^2_\mathrm{decoupled}
  < \|\nabla_\theta \mathcal{L}_u + \nabla_\theta \mathcal{L}_g\|^2_\mathrm{shared}
\]
whenever the encoder Jacobians for understanding and generation are
non-parallel --- which holds whenever the two tasks require different
visual features, a natural assumption.

\section{Experimental Findings}

\textbf{Understanding benchmarks (VQA, image classification):}
The task-decoupled model achieves 3.2\% average relative improvement
over the best single-task baseline and 5.8\% over the naive joint model.

\textbf{Generation benchmarks (FID on COCO, 256$\times$256):}

\begin{center}
\begin{tabular}{lcc}
\toprule
Model & FID$\downarrow$ & Params \\
\midrule
Generation-only specialist & 7.4 & 1B \\
Naive joint UMM & 10.3 & 1B \\
\textbf{Task-decoupled (ours)} & \textbf{8.1} & 1.1B \\
\bottomrule
\end{tabular}
\end{center}

\textbf{Parameter overhead:} The two specialised encoders add 12\%
extra parameters vs.\ the fully shared model, but outperform it on both
tasks simultaneously.

\section{Ablations and Interpretation}

\textbf{Which layers to specialise?} Specialising only the visual
tokeniser (earliest layers) provides the largest gain. Specialising
all layers including late transformer blocks hurts by breaking shared
semantic grounding.

\textbf{Loss weight sensitivity ($\lambda_u, \lambda_g$):} The naive
joint model is highly sensitive to loss weighting: accuracy fluctuates
by up to 8\% as $\lambda_u / \lambda_g$ varies from 0.5 to 2.0.
The decoupled model is robust across this range.

\textbf{Scaling (300M, 1B, 3B):} The performance gains persist at all
three scales, suggesting the conflict is a structural, not a capacity,
issue.

\textbf{CKA trajectories:} During training, understanding and generation
encoder representations diverge (lower cross-task CKA), while semantic
backbone representations converge (higher cross-task CKA) --- matching
the specialisation-then-fusion theory.

\section*{Reference}

Penghao Wu, Haiwen Diao, Weichen Fan, Lewei Lu, Dahua Lin, and Ziwei Liu.
\textbf{Uncovering Understanding-Generation Synergy in Native Unified
Multimodal Models: From Representation, Task to System.}
arXiv:2609.01607, September 2026.
\url{https://arxiv.org/abs/2609.01607}

\end{document}
